\section{The Lean check: process record}\label{supp:lean}

The Lean file was written by the main assistant. Lean could not be installed in the assistant's cloud workspace,
whose network policy blocks the servers that distribute Lean, so the file was compiled on the author's computer.
The author installed Lean and created the project; the assistant read the build output from the project folder
through the link between the Claude desktop application and the author's computer, wrote corrected versions of the
file into the same folder, and the author ran the next build (Figure~\ref{fig:leanloop}).

\begin{figure}[ht]
\centering
\begin{tikzpicture}[node distance=6mm and 14mm,
  box/.style={draw,rounded corners,align=center,font=\small,minimum height=10mm,text width=34mm},
  >={Stealth[length=2mm]}]
\node[box] (a) {Assistant writes or\\corrects \texttt{Basic.lean}};
\node[box,right=of a] (b) {File written into the\\author's project folder};
\node[box,below=of b] (c) {Author runs\\\texttt{lake build}};
\node[box,left=of c] (d) {Assistant reads\\\texttt{build\_log.txt}};
\draw[->] (a) -- (b);
\draw[->] (b) -- (c);
\draw[->] (c) -- (d);
\draw[->] (d) -- node[left,font=\footnotesize,align=center]{errors}(a);
\node[box,right=of c,text width=30mm] (e) {Build succeeds;\\axiom report};
\draw[->] (c) -- (e);
\end{tikzpicture}
\caption{The build-and-repair loop used for the Lean check.}\label{fig:leanloop}
\end{figure}

Table~\ref{tab:leantime} gives the timeline. Times are in UTC and are taken from the session records and from the
modification times of the files in the project folder.

\begin{table}[ht]
\centering\small
\begin{tabular}{@{}>{\raggedright\arraybackslash}p{3.3cm}>{\raggedright\arraybackslash}p{11.0cm}@{}}
\toprule
\textbf{Time (UTC)} & \textbf{Step}\\
\midrule
21 Oct, 21:38 & The assistant writes a first Lean file for the results in Table~\ref{tab:lean}; it cannot be compiled
  in the cloud workspace.\\
24 Oct, 17:08 & The author asks how a formal proof check can be used; the assistant describes what is feasible
  (the loss-ratio steps), what is hard (the representation step) and how to install Lean.\\
24 Oct, 17:26--17:31 & The author installs the Lean extension for the code editor and the Lean toolchain manager.\\
24 Oct, 17:35--17:37 & The author creates the project \texttt{loss\_ratio\_lean} with Mathlib; Lean and Mathlib
  version~4.36.0-rc1; the precompiled Mathlib files are downloaded.\\
24 Oct, 17:38--17:39 & The assistant inspects the project configuration, confirms that it is complete, and places
  the Lean file in it, replacing the template placeholder.\\
25 Oct, 12:47--12:53 & First attempt to build: the command \texttt{lake} is not found, because the toolchain had not
  been added to the terminal's search path; fixed by loading the toolchain environment.\\
25 Oct, 12:55 & Build 1: four errors, none in a definition or statement. One proof stopped one algebraic step short;
  in three proofs Lean's simplifier rewrote $\mathrm{const}(-c)$ as $-\mathrm{const}(c)$, so that cash additivity
  no longer applied.\\
25 Oct, 12:57 & The assistant writes a corrected file and appends \texttt{\#print axioms} for ten main theorems.\\
25 Oct, 13:04 & Build 2: one error, in the converse construction, where the goal was displayed in an unexpanded
  form. The axiom report shows the three standard axioms for nine theorems and \texttt{sorryAx} for the converse.\\
25 Oct, 13:05 & The assistant adds one line that states the goal explicitly.\\
25 Oct, 13:08 & Build 3 succeeds. All ten theorems depend only on the three standard axioms.\\
25 Oct, 13:09 & The assistant compares the Lean definitions of \axP, \axN, \axCM{} and \axLI{} and the theorem
  statements with Section~\ref{sec:axioms}.\\
25 Oct, 13:37 & At the author's request, the project's own build output is deleted and the file is rebuilt from
  scratch: success in 8.8\,s, with the same axiom report.\\
\bottomrule
\end{tabular}
\caption{Timeline of the Lean check.}\label{tab:leantime}
\end{table}

